% ====================================================================
% ФАЙЛ: tasks/02_mkt/mkt_heat_and_phase_transitions_bank_part1.tex
% ТЕМА: КОЛИЧЕСТВО ТЕПЛОТЫ. УДЕЛЬНАЯ ТЕПЛОЕМКОСТЬ. ФАЗОВЫЕ ПЕРЕХОДЫ
% ПАЧКА 1: ВАРИАНТЫ 1–5
% ====================================================================

% === ЗАДАЧА 1 ===
% Тег: #МКТ #ТЕПЛОТА #КИМ_08 #ВАРИАНТ_03
\begin{taskbasic}[code={\label{task:mkt_heat_v3_n8}}]
\textbf{Задача 1.} На рисунке показан график зависимости температуры вещества в процессе теплообмена с окружающей средой от отдаваемого им количества теплоты. Масса вещества равна $150$~г. Первоначально вещество было в газообразном состоянии. Какова удельная теплота парообразования вещества?

\nopagebreak\smallskip
\begin{center}
\begin{tikzpicture}[scale=1.1, every node/.style={font=\footnotesize}]
    \draw[xstep=1, ystep=1, gray!30, thin] (0,0) grid (6,4);
    \draw[-{Stealth[scale=0.8]}, black, thick] (0,0) -- (6.5,0) node[right] {$Q, 10^5\text{ Дж}$};
    \draw[-{Stealth[scale=0.8]}, black, thick] (0,0) -- (0,4.5) node[above] {$t, \,^{\circ}\text{C}$};
    
    \coordinate (A) at (0,4);
    \coordinate (B) at (1,2.5);
    \coordinate (C) at (4,2.5);
    \coordinate (D) at (6,0.5);
    
    \draw[ultra thick, brandPrimary] (A) -- (B) -- (C) -- (D);
    
    \draw[dashed, gray] (1,0) -- (1,2.5);
    \draw[dashed, gray] (2,0) -- (2,2.5);
    \draw[dashed, gray] (3,0) -- (3,2.5);
    \draw[dashed, gray] (4,0) -- (4,2.5);
    \draw[dashed, gray] (5,0) -- (5,1.5);
    \draw[dashed, gray] (6,0) -- (6,0.5);
    
    \node[below] at (1,0) {1};
    \node[below] at (2,0) {2};
    \node[below] at (3,0) {3};
    \node[below] at (4,0) {4};
    \node[below] at (5,0) {5};
    \node[below] at (6,0) {6};
    \node[below left] at (0,0) {0};
\end{tikzpicture}
\end{center}

\kim{8}
\end{taskbasic}
\answer{2000}

% === ЗАДАЧА 2 ===
% Тег: #МКТ #ТЕПЛОТА #КИМ_08 #ВАРИАНТ_04
\begin{taskbasic}[code={\label{task:mkt_heat_v4_n8}}]
\textbf{Задача 2.} На рисунке показан график зависимости температуры вещества в процессе теплообмена с окружающей средой от отдаваемого им количества теплоты. Масса вещества равна $1{,}2$~кг. Первоначально вещество было в жидком состоянии. Какова удельная теплота плавления вещества?

\nopagebreak\smallskip
\begin{center}
\begin{tikzpicture}[scale=1.1, every node/.style={font=\footnotesize}]
    \draw[xstep=1, ystep=1, gray!30, thin] (0,0) grid (6,4);
    \draw[-{Stealth[scale=0.8]}, black, thick] (0,0) -- (6.5,0) node[right] {$Q, 10^5\text{ Дж}$};
    \draw[-{Stealth[scale=0.8]}, black, thick] (0,0) -- (0,4.5) node[above] {$t, \,^{\circ}\text{C}$};
    
    \coordinate (A) at (0,4);
    \coordinate (B) at (1,2.5);
    \coordinate (C) at (4,2.5);
    \coordinate (D) at (6,0.5);
    
    \draw[ultra thick, brandPrimary] (A) -- (B) -- (C) -- (D);
    
    \draw[dashed, gray] (1,0) -- (1,2.5);
    \draw[dashed, gray] (2,0) -- (2,2.5);
    \draw[dashed, gray] (3,0) -- (3,2.5);
    \draw[dashed, gray] (4,0) -- (4,2.5);
    \draw[dashed, gray] (5,0) -- (5,1.5);
    \draw[dashed, gray] (6,0) -- (6,0.5);
    
    \node[below] at (1,0) {1};
    \node[below] at (2,0) {2};
    \node[below] at (3,0) {3};
    \node[below] at (4,0) {4};
    \node[below] at (5,0) {5};
    \node[below] at (6,0) {6};
    \node[below left] at (0,0) {0};
\end{tikzpicture}
\end{center}

\kim{8}
\end{taskbasic}
\answer{250}


% ====================================================================
% ФАЙЛ: tasks/02_mkt/mkt_heat_and_phase_transitions_bank_part2.tex
% ТЕМА: КОЛИЧЕСТВО ТЕПЛОТЫ. УДЕЛЬНАЯ ТЕПЛОЕМКОСТЬ. ФАЗОВЫЕ ПЕРЕХОДЫ
% ПАЧКА 2: ВАРИАНТЫ 6–15
% ====================================================================



% === ЗАДАЧА 2 ===
% Тег: #МКТ #ТЕПЛОТА #КИМ_08 #ВАРИАНТ_07
\begin{taskbasic}[code={\label{task:mkt_heat_v7_n8}}]
\textbf{Задача 2.} Твёрдое тело массой $4$~кг нагревают. На рисунке представлен график зависимости абсолютной температуры тела от полученного им количества теплоты. Чему равна удельная теплоёмкость вещества, из которого состоит тело?

\nopagebreak\smallskip
\begin{center}
\begin{tikzpicture}[scale=1.1, every node/.style={font=\footnotesize}]
    \draw[xstep=1, ystep=1, gray!30, thin] (0,0) grid (4,4);
    \draw[-{Stealth[scale=0.8]}, black, thick] (0,0) -- (4.5,0) node[right] {$Q, \text{ кДж}$};
    \draw[-{Stealth[scale=0.8]}, black, thick] (0,0) -- (0,4.5) node[above] {$t, \,^{\circ}\text{C}$};
    
    \coordinate (A) at (0,1);
    \coordinate (B) at (4,3.4);
    
    \draw[ultra thick, brandPrimary] (A) -- (B);
    
    \draw[dashed, gray] (1,0) -- (1,1.6);
    \draw[dashed, gray] (2,0) -- (2,2.2);
    \draw[dashed, gray] (3,0) -- (3,2.8);
    \draw[dashed, gray] (4,0) -- (4,3.4);
    
    \node[left] at (0,1) {100};
    \node[left] at (0,2) {150};
    \node[left] at (0,3) {200};
    \node[below] at (1,0) {120};
    \node[below] at (2,0) {240};
    \node[below] at (3,0) {360};
    \node[below left] at (0,0) {0};
\end{tikzpicture}
\end{center}

\kim{8}
\end{taskbasic}
\answer{150}

% === ЗАДАЧА 3 ===
% Тег: #МКТ #ТЕПЛОТА #КИМ_08 #ВАРИАНТ_08
\begin{taskbasic}[code={\label{task:mkt_heat_v8_n8}}]
\textbf{Задача 3.} Твёрдое тело нагревают. На рисунке представлен график зависимости абсолютной температуры тела от полученного им количества теплоты. Удельная теплоёмкость вещества, из которого состоит тело, равна $720$~Дж/(кг$\cdot$К). Чему равна масса тела?

\nopagebreak\smallskip
\begin{center}
\begin{tikzpicture}[scale=1.1, every node/.style={font=\footnotesize}]
    \draw[xstep=1, ystep=1, gray!30, thin] (0,0) grid (4,4);
    \draw[-{Stealth[scale=0.8]}, black, thick] (0,0) -- (4.5,0) node[right] {$Q, \text{ кДж}$};
    \draw[-{Stealth[scale=0.8]}, black, thick] (0,0) -- (0,4.5) node[above] {$t, \,^{\circ}\text{C}$};
    
    \coordinate (A) at (0,1);
    \coordinate (B) at (4,3.4);
    
    \draw[ultra thick, brandPrimary] (A) -- (B);
    
    \draw[dashed, gray] (1,0) -- (1,1.6);
    \draw[dashed, gray] (2,0) -- (2,2.2);
    \draw[dashed, gray] (3,0) -- (3,2.8);
    \draw[dashed, gray] (4,0) -- (4,3.4);
    
    \node[left] at (0,1) {100};
    \node[left] at (0,2) {150};
    \node[left] at (0,3) {200};
    \node[below] at (1,0) {120};
    \node[below] at (2,0) {240};
    \node[below] at (3,0) {360};
    \node[below left] at (0,0) {0};
\end{tikzpicture}
\end{center}

\kim{8}
\end{taskbasic}
\answer{3}